\chapter{Manual de usuario}
\label{chap:user-manual}

\section{Requisitos previos}
\section{Instalación y puesta en marcha}

\section{Acceso a la aplicación}
\label{sec:user-manual-access}

La aplicación dispone de un sistema de autenticación mediante correo
electrónico y contraseña. Desde la pantalla de acceso, el usuario puede iniciar
sesión con una cuenta existente o acceder al formulario de registro mediante el
enlace \emph{Regístrate}.

\subsection{Creación de una cuenta}

Para crear una cuenta nueva, el usuario debe completar los campos de nombre,
apellido, correo electrónico, contraseña y nombre de la empresa, tal y como se
muestra en la Figura~\ref{fig:user-manual-register}. Todos los campos son
obligatorios, el correo debe tener un formato válido y la contraseña debe
contener al menos ocho caracteres. Una vez cumplimentado el formulario, se debe
pulsar el botón \emph{Crear cuenta}.

\begin{figure}[H]
  \centering
  \includegraphics[width=0.72\textwidth,trim=250bp 45bp 250bp 0,clip]{screenshots/auth-register.png}
  \caption{Formulario de creación de una cuenta en el prototipo actual.}
  \label{fig:user-manual-register}
\end{figure}

Si ya se dispone de una cuenta, el enlace \emph{Inicia sesión}, situado bajo el
formulario, permite regresar a la pantalla de acceso.

\subsection{Inicio de sesión}

En la pantalla mostrada en la Figura~\ref{fig:user-manual-login}, el usuario
debe introducir su correo electrónico y su contraseña y pulsar
\emph{Iniciar sesión}. El sistema comprueba que el correo tenga un formato
válido y que la contraseña no esté vacía antes de enviar la solicitud.

\begin{figure}[H]
  \centering
  \includegraphics[width=0.72\textwidth,trim=250bp 45bp 250bp 0,clip]{screenshots/auth-login.png}
  \caption{Pantalla de inicio de sesión del prototipo actual.}
  \label{fig:user-manual-login}
\end{figure}

Si las credenciales no son correctas, la aplicación informa al usuario mediante
un mensaje de error. Cuando la autenticación se completa correctamente, el
usuario es redirigido a la página solicitada inicialmente o, si no existe una
redirección pendiente, a la pantalla principal.

\section{Gestión de oportunidades}
\section{Configuración de pipelines y workflows}
\section{Gestión del CRM}
